\begin{solapartat}
\punts{0,25} Flux inicial:
\[ \phi_m = \vec{B}\cdot\vec{S} = B\cdot S = B\cdot b\cdot x = 0,4\cdot 0,04\cdot 0,03 = 4,80\times 10^{-4}\ \mathrm{Wb} \]

\punts{0,5} Tal com es pot veure a la figura, el corrent induït en el circuit gira en sentit horari. Aplicant la regla de la mà dreta, veiem que aquest corrent induït genera un camp magnètic induït cap a dins del paper (és a dir, té el sentit contrari que el camp magnètic extern).
\figura[0.4]{sol_fig1.png}
A partir de la llei de Lenz, sabem que el corrent induït crea un camp magnètic induït que s'oposa a la variació de flux magnètic.

Per tant, podem deduir que el moviment de la vareta ha de ser tal que el flux magnètic degut al camp magnètic extern augmenti. Això implica que la vareta s'està movent cap a la dreta:
\figura[0.33]{sol_fig2.png}

\punts{0,5} Flux magnètic:
\[ \phi_m = \vec{B}\cdot\vec{S} = B\cdot S = B\cdot b\cdot x = B\cdot b\cdot (x_0 + vt) \]
FEM induïda:
\[ \varepsilon = \left|\frac{d\phi_m}{dt}\right| = B\cdot b\cdot v = 4,80\cdot 10^{-2}\ \mathrm{V} \]
\end{solapartat}

\begin{solapartat}
\punts{0,5} La llei de Lorentz indica que apareix una força magnètica sobre les càrregues en moviment, perpendicular a la velocitat de les càrregues i al camp magnètic.

En el cas dels fils conductors, la força sobre un fil de longitud $l$ pel qual circula una intensitat $I$ és la següent:
\[ \vec{F}_m = I(\vec{l}\times\vec{B}) \]
En el cas descrit en l'enunciat, la vareta i el camp magnètic són perpendiculars, com s'indica a la figura. Per tant, apareixerà una força magnètica sobre la vareta.

\punts{0,5} Com que són perpendiculars, el mòdul de la força s'expressa així:
\[ F_m = I\cdot l\cdot B\cdot \sin(90) = 6\cdot 10^{-3}\cdot 0,04\cdot 0,4\cdot 1 = 9,60\cdot 10^{-5}\ \mathrm{N} \]
\punts{0,25} Dibuix de la força magnètica, perpendicular a $\vec{l}$ i $\vec{B}$ i segons la regla de la mà dreta en el sentit indicat a la figura:
\figura[0.3]{sol_fig3.png}
També es pot justificar el sentit de la força a partir de la llei de Lenz. Com que s'oposa al canvi i la vareta es desplaça cap a la dreta, la força tindrà el sentit oposat, és a dir, cap a l'esquerra.
\end{solapartat}
